\section{Introduction}
\label{sec:introduction}

Generation constructs a sample from an initial state; repair resumes that construction after the sample has been disturbed. The term \emph{regeneration} comes here from the biological motivation of neural cellular automata. Growing NCA distinguishes growth from a seed, persistence of the resulting pattern, and recovery after damage~\cite{nca}. We adopt regeneration in this behavioral sense, rather than as a separate category of probabilistic model or a capability exclusive to cellular automata.

Neural Cellular Automata Manifold (NCAM) supplies the autoencoding lineage: an encoder selects a cellular growth program capable of reconstructing a target image~\cite{ncam}. Regenerative Autoencoders (ReGAE) retain the idea of a reusable reconstruction process while replacing the cellular decoder with a shape-preserving denoiser. The denoiser acts on a corrupted version of the desired state and directly predicts a corrected state. Its output becomes the input to the next update. Starting with noise gives generation; starting with a disturbed reconstruction gives regeneration. In the conditional form, a latent code identifies the intended reconstruction throughout either process.

Neither an iteration index nor a corruption-strength label is supplied to the ReGAE denoiser. At fixed conditioning, the same state is therefore handled by the same learned map wherever it occurs in a trajectory. This motivates viewing generation and repair as movement under one correction field, learned across a range of corruptions. The interpretation is related to diffusion and flow models, but their prediction targets, dependence on time, and sampling rules need to be separated carefully. A single map does not by itself establish a scalar optimization landscape or a globally attracting solution.

\begin{figure*}[t]
  \centering
  % TikZ fragment for a full-width figure. Requires tikz and arrows.meta.
\begin{tikzpicture}[
  x=1cm,y=1cm,
  font=\sffamily\footnotesize,
  >={Latex[length=2mm,width=1.35mm]},
  box/.style={draw=black!45,rounded corners=2pt,line width=.55pt,
              fill=black!3,align=center,minimum height=.82cm,inner sep=3pt},
  net/.style={box,draw=blue!55!black,fill=blue!7},
  flow/.style={->,line width=.65pt,draw=black!70},
  note/.style={font=\sffamily\scriptsize,text=black!70,align=center}
]
% Pixel-level reconstruction.
\node[anchor=west,font=\sffamily\bfseries\small] at (-.1,1.55)
  {Stage 1: patch autoencoding};
\node[box,text width=2.1cm] (rgb) at (1.05,.65)
  {RGB patch $x$\\$32\!\times\!32\!\times\!3$};
\node[net,text width=1.4cm] (enc) at (3.95,.65)
  {Encoder\\$E_{\phi}$};
\node[box,text width=1.9cm] (code) at (6.4,.65)
  {Patch code $z$\\256 values};
\node[net,text width=2.0cm] (dec) at (10.1,.65)
  {Conditional\\denoiser $D_{\theta}$};
\node[box,text width=2.45cm] (recon) at (13.65,.65)
  {Reconstructed patch\\$\widehat{x}$};
\node[box,text width=2.1cm] (corruptx) at (6.4,-.53)
  {Corrupted pixels\\$\widetilde{x}$};
\draw[flow] (rgb)--(enc);
\draw[flow] (enc)--(code);
\draw[flow] (code)--node[note,above]{conditioning}(dec);
\draw[flow] (dec)--(recon);
\draw[flow] (rgb.south) |- node[note,pos=.70,above]{corruption} (corruptx.west);
\draw[flow] (corruptx.east) -| (dec.south);
\draw[black!15,line width=.5pt] (-.1,-1.12)--(15.0,-1.12);
% Spatial latent prediction, followed by patchwise reconstruction.
\node[anchor=west,font=\sffamily\bfseries\small] at (-.1,-1.57)
  {Stage 2: spatial latent prediction};
\node[box,text width=2.1cm] (window) at (1.05,-2.44)
  {Encoded window $z^-$\\$32\!\times\!32\!\times\!256$};
\node[box,text width=1.65cm] (corruptz) at (4.05,-2.44)
  {Corruption\\$\widetilde{z}^-$};
\node[net,text width=1.8cm] (gen) at (7.05,-2.44)
  {Gencoder\\$G_{\psi}$};
\node[box,text width=2.45cm] (shifted) at (10.35,-2.44)
  {Predicted shifted\\latent window};
\node[box,text width=2.45cm] (image) at (13.65,-2.44)
  {Decode each code\\$1024\!\times\!1024$ pixels};
\draw[flow] (window)--(corruptz);
\draw[flow] (corruptz)--(gen);
\draw[flow] (gen)--(shifted);
\draw[flow] (shifted)--(image);
\node[note] at (7.55,-3.25)
  {$D_{\theta}$ and $G_{\psi}$ reuse the same gated denoiser family with separate parameters.};
\end{tikzpicture}

  \caption{Conditional and unconditional members of the ReGAE family. Both denoisers return a state with the same shape as their input, allowing the output to be fed back for successive updates. Generation initializes the state with noise; regeneration re-enters the same loop with a perturbed state. In the conditional model, an encoder supplies a latent tensor that is retained during regeneration. The unconditional model has no encoder-derived conditioning and can model data or learned representations. The two rows illustrate variants, not mandatory levels of one hierarchy.}
  \label{fig:architecture}
\end{figure*}

The broader question motivating this work is what it takes to learn, or \emph{memorize}, a dataset. Calling a model a memorizer does not explain how that memory is organized, which properties it retains, or how it makes them available for generation. A dataset contains information about individual samples and, potentially, relationships among them. Learning to reconstruct a sample addresses one part of this problem. Learning the distribution or organization of its representations addresses another. These are useful questions even before asking whether a system extrapolates beyond its training data.

ReGAE provides a reusable modeling component for investigating those questions. An encoder can concentrate information about a particular sample into a code, while a conditional denoiser learns how to recover the sample from different corrupted states. Removing the encoder and its conditioning pathway gives an unconditional member of the family: a denoiser whose evolving state supplies its context. Such a model can operate directly on data or on representations produced by another model. Conditional and unconditional ReGAE instances can therefore be composed, without requiring a particular number of levels or a particular kind of dataset.

The experiments presented here use a large image of Pandora's Cluster as an example dataset. An image-patch autoencoder learns local appearance, while a second, unconditional ReGAE learns the spatial organization of the resulting patch codes. We call this latent-space instance the \emph{Gencoder}. Predicting codes for a slightly displaced view couples regeneration to a spatial transition. This arrangement is appropriate because spatial relationships are an important part of the information in this dataset. It is an application of the family, rather than its definition.

A different dataset would suggest a different decomposition. For a collection such as CIFAR-10, the first autoencoder could represent each complete image. A subsequent model could learn the distribution of image codes or, with suitable inputs and supervision, relationships between image and class representations. There would be no reason to impose the spatial arrangement used for Pandora's Cluster. This example illustrates the intended scope; it is not an additional experiment reported here.

The defining design choice is thus an iteratively reusable, shape-preserving denoiser, with optional latent conditioning. Its internal architecture, data dimensionality, reconstruction loss, and placement within a hierarchy are configurable. We describe this general construction first, then use the astronomical-image proof of concept to examine the distinction between sample reconstruction and modeling relationships among samples. The discussion connects regeneration to denoising dynamics, architectural reuse, and compression, while keeping the broader proposal distinct from the evidence supplied by the existing experiments.
